\SourceLeaf{91}%
Quarta cum secunda fit identica. Adiumento harum aequationum tres incognitarum eliminare licet, quo pacto omnes sedecim iam ad duas reductae sunt.

\Article{17}
Ut vero determinationem completam nanciscamur, investigare conveniet multitudinem solutionum congruentiae
\[
1+\alpha+\beta+\gamma\equiv0\ (\mathrm{mod.}\ p)
\]
designantibus $\alpha,\beta,\gamma$ indefinite numeros e complexibus $A,B,C$. Manifesto valor $\alpha=p-1$ non est admissibilis, quum fieri nequeat $\beta+\gamma\equiv0$: substituendo itaque pro $\alpha$ deinceps valores reliquos, prodibunt $h,i,k,l$ valores ipsius $1+\alpha$ ad $A,B,C,D$ resp. pertinentes. Pro quovis autem valore \emph{dato} ipsius $1+\alpha$ ad $A$ pertinente, puta pro $1+\alpha=\alpha^0$, congruentia $\alpha^0+\beta+\gamma\equiv0$ totidem solutiones admittet, quot congruentia $1+\beta'+\gamma'\equiv0$ (statuendo scilicet $\beta\equiv\alpha^0\beta'$, $\gamma\equiv\alpha^0\gamma'$), i. e. solutiones $(12)=m$. Perinde pro quovis valore dato ipsius $1+\alpha$ ad $B$ pertinente, puta pro $1+\alpha=\beta^0$, congruentia $\beta^0+\beta+\gamma\equiv0$ totidem solutiones habebit, quot haec $1+\alpha'+\beta'\equiv0$ (scilicet statuendo $\beta\equiv\beta^0\alpha'$, $\gamma\equiv\beta^0\beta'$), i. e. solutiones $(01)=i$. Similiter pro quolibet valore dato ipsius $1+\alpha$ ad $C$ pertinente, puta pro $1+\alpha=\gamma^0$, congruentia $\gamma^0+\beta+\gamma\equiv0$ totidem modis diversis solvi poterit, quot haec $1+\delta+\alpha'\equiv0$ (nempe statuendo $\beta\equiv\gamma^0\delta$, $\gamma\equiv\gamma^0\alpha'$), i. e. solutionum multitudo erit $(03)=l$. Denique pro quovis valore dato ipsius $1+\alpha$ ad $D$ pertinente, puta pro $1+\alpha=\delta^0$, congruentia $\delta^0+\beta+\gamma\equiv0$ totidem solutiones habebit, quot haec $1+\gamma'+\delta'\equiv0$ (statuendo $\beta\equiv\delta^0\gamma'$, $\gamma\equiv\delta^0\delta'$), i. e. $(23)=m$ solutiones. Omnibus itaque collectis, patet, congruentiam $1+\alpha+\beta+\gamma\equiv0$ admittere
\[
hm+ii+kl+lm
\]
solutiones diversas.

Prorsus vero simili modo eruimus, si pro $\beta$ singuli deinceps numeri complexus $B$ substituantur, summam $1+\beta$ obtinere resp. $(10),(11),(12),(13)$ sive $i,l,m,m$ valores ad $A,B,C,D$ pertinentes, et pro quovis valore \emph{dato} ipsius $1+\beta$ ad hos complexus pertinente, congruentiam $1+\beta+\alpha+\gamma\equiv0$ resp. $(02),(31),(20),(13)$ sive $k,m,k,m$ solutiones diversas admittere, ita ut multitudo omnium solutionum fiat

\SourceLeaf{92}%
\[
=ik+lm+km+mm
\]
Ad eundem valorem perducimur, si evolutionem considerationi valorum summae $1+\gamma$ superstruimus.

\Article{18}
Ex hac duplici eiusdem multitudinis expressione nanciscimur aequationem:
\[
0=hm+ii+kl-ik-km-mm
\]
atque hinc, eliminando $h$ adiumento aequationis $h=2m-k-1$,
\[
0=(k-m)^2+ii+kl-ik-kk-m
\]
Sed duae aequationes ultimae art. 16 suppeditant $k=\frac12(l+i)$, quo valore substituto $ii+kl-ik-kk$ transit in $\frac14(l-i)^2$, adeoque aequatio praecedens, per $4$ multiplicata, in hanc
\[
0=4(k-m)^2+(l-i)^2-4m
\]
Hinc, quoniam $4m=2(k+m)-2(k-m)=2n-2(k-m)$, sequitur
\[
2n=4(k-m)^2+2(k-m)+(l-i)^2
\]
sive
\[
8n+1=(4(k-m)+1)^2+4(l-i)^2
\]
Statuendo itaque
\[
4(k-m)+1=a,\quad2l-2i=b
\]
habebimus
\[
p=aa+bb
\]

Sed constat, $p$ unico tantum modo in duo quadrata discerpi posse, quorum alterum impar accipi debet pro $aa$, alterum par pro $bb$, ita ut $aa,bb$ sint numeri ex asse determinati. Sed etiam $a$ ipse erit numerus prorsus determinatus: radix enim quadrati positive accipi debet. vel negative, prout radix positiva est formae $4M+1$ vel $4M+3$. De determinatione signi ipsius $b$ mox loquemur.

Iam combinatis his novis aequationibus cum tribus ultimis art. 16, quinque numeri $h,i,k,l,m$ per $a,b$ et $n$ penitus determinantur sequenti modo:

\SourceLeaf{93}%
\[
\begin{aligned}
8h&=4n-3a-5\\
8i&=4n+a-2b-1\\
8k&=4n+a-1\\
8l&=4n+a+2b-1\\
8m&=4n-a+1
\end{aligned}
\]

Si loco ipsius $n$ modulum $p$ introducere malumus, schema $S$, singulis terminis ad evitandas fractiones per $16$ multiplicatis, ita se habet:
\[
\begin{array}{l|l|l|l}
p-6a-11&p+2a-4b-3&p+2a-3&p+2a+4b-3\\
p+2a-4b-3&p+2a+4b-3&p-2a+1&p-2a+1\\
p+2a-3&p-2a+1&p+2a-3&p-2a+1\\
p+2a+4b-3&p-2a+1&p-2a+1&p+2a-4b-3
\end{array}
\]

\Article{19}
Superest, ut signum ipsi $b$ tribuendum assignare doceamus. Iam supra, art. 10, monuimus, distinctionem inter complexus $B$ et $D$, per se non essentialem, ab electione numeri $f$ pendere, pro quo alterutra radix congruentiae $xx\equiv-1$ accipi debet, illasque inter se permutari, si loco alterius radicis altera adoptetur. Iam quum inspectio schematis modo allati doceat, similem permutationem cum mutatione signi ipsius $b$ cohaerere, praevidere licet, nexum inter signum ipsius $b$ atque numerum $f$ exstare debere. Quem ut cognoscamus, ante omnia observamus, si, denotante $\mu$ integrum non negativum, pro $z$ accipiantur omnes numeri $1,2,3\dotsfour p-1$, fieri secundum modulum $p$, vel $\Sigma z^\mu\equiv0$, vel $\Sigma z^\mu\equiv-1$, prout $\mu$ vel non-divisibilis sit per $p-1$, vel divisibilis. Pars posterior theorematis inde patet, quod pro valore ipsius $\mu$ per $p-1$ divisibili, habetur $z^\mu\equiv1$: partem priorem vero ita demonstramus. Denotante $g$ radicem primitivam, omnes $z$ convenient cum residuis minimis omnium $g^y$, accipiendo pro $y$ omnes numeros $0,1,2,3\dotsfour p-2$, eritque adeo $\Sigma z^\mu\equiv\Sigma g^{\mu\pdot y}$. Sed fit
\[
\begin{gathered}
\Sigma g^{\mu y}=\frac{g^{\mu(p-1)}-1}{g^\mu-1},\quad\text{adeoque}\\
(g^\mu-1)\Sigma z^\mu\equiv g^{\mu(p-1)}-1\equiv0
\end{gathered}
\]
Hinc vero sequitur, quoniam pro valore ipsius $\mu$ per $p-1$ non-divisibili $g^\mu$ ipsi $1$ congruus sive $g^\mu-1$ per $p$ divisibilis esse nequit, $\Sigma z^\mu\equiv0$. Q. E. D.

\SourceLeaf{94}%
Iam si potestas $(z^4+1)^{\frac14(p-1)}$ secundum theorema binomiale evolvitur, per lemma praec. fiet
\[
\Sigma(z^4+1)^{\frac14(p-1)}\equiv-2\ (\mathrm{mod.}\ p)
\]
Sed residua minima omnium $z^4$ exhibent omnes numeros $A$, quovis quater occurrente; habebimus itaque inter residua minima ipsius $z^4+1$
\[
\begin{array}{rl}
4(00)&\text{ad }A\\
4(01)&\text{ad }B\\
4(02)&\text{ad }C\\
4(03)&\text{ad }D
\end{array}
\]
pertinentia, quatuorque erunt $=0$ (puta pro $z^4\equiv p-1$). Hinc, considerando criteria complexuum $A,B,C,D$, deducimus
\[
\begin{aligned}
\Sigma(z^4+1)^{\frac14(p-1)}\equiv{}&4(00)+4f(01)\\
&-4(02)-4f(03)
\end{aligned}
\]
adeoque
\[
-2\equiv4(00)+4f(01)-4(02)-4f(03)
\]
sive substitutis pro $(00),(01)$ etc. valoribus in art. praec. inventis,
\[
-2\equiv-2a-2-2bf
\]
Hinc itaque colligimus, semper fieri debere $a+bf\equiv0$, sive, multiplicando per $f$,
\[
b\equiv af
\]
quae congruentia determinationi signi ipsius $b$, si numerus $f$ iam electus est, vel determinationi numeri $f$, si signum ipsius $b$ aliunde praescribitur, inservit.

\Article{20}
Postquam problema nostrum pro modulis formae $8n+1$ complete solvimus, progredimur ad casum alterum, ubi $p$ est formae $8n+5$: quem eo brevius absolvere licebit, quod omnia ratiocinia parum a praecedentibus differunt.

Quum pro tali modulo $-1$ ad classem $C$ pertineat, complementa numerorum complexuum $A,B,C,D$ ad summam $p$, in classibus $C,D,A,B$ resp. contenta erunt. Hinc facile colligitur

\SourceLeaf{95}%
\begin{center}
\begin{tabular}{c|l}
signum&\shortstack{denotare multitudinem\\solutionum congruentiae}\\\hline
$(00)$&$1+\alpha+\gamma\equiv0$\\
$(01)$&$1+\alpha+\delta\equiv0$\\
$(02)$&$1+\alpha+\alpha'\equiv0$\\
$(03)$&$1+\alpha+\beta\equiv0$\\
$(10)$&$1+\beta+\gamma\equiv0$\\
$(11)$&$1+\beta+\delta\equiv0$\\
$(12)$&$1+\beta+\alpha\equiv0$\\
$(13)$&$1+\beta+\beta'\equiv0$\\
$(20)$&$1+\gamma+\gamma'\equiv0$\\
$(21)$&$1+\gamma+\delta\equiv0$\\
$(22)$&$1+\gamma+\alpha\equiv0$\\
$(23)$&$1+\gamma+\beta\equiv0$\\
$(30)$&$1+\delta+\gamma\equiv0$\\
$(31)$&$1+\delta+\delta'\equiv0$\\
$(32)$&$1+\delta+\alpha\equiv0$\\
$(33)$&$1+\delta+\beta\equiv0$
\end{tabular}
\end{center}
unde statim habentur sex aequationes:
\[
\begin{gathered}
(00)=(22),\quad(01)=(32),\quad(03)=(12),\\
(10)=(23),\quad(11)=(33),\quad(21)=(30)
\end{gathered}
\]

Multiplicando congruentiam $1+\alpha+\gamma\equiv0$ per numerum $\gamma'$ e complexu $C$ ita electum, ut fiat $\gamma\gamma'\equiv1$, accipiendoque pro $\gamma''$ residuum minimum producti $\alpha\gamma'$, quod manifesto quoque complexui $C$ adnumerandum erit, prodit $\gamma'+\gamma''+1\equiv0$, unde colligimus $(00)=(20)$.

Prorsus simili modo habentur aequationes $(01)=(13)$, $(03)=(31)$, $(10)=(11)=(21)$.

Adiumento harum undecim aequationum sedecim incognitas nostras ad quinque reducere, schemaque $S$ ita exhibere possumus:
\[
\begin{array}{llll}
h,&i,&k,&l\\
m,&m,&l,&i\\
h,&m,&h,&m\\
m,&l,&i,&m
\end{array}
\]

